\documentclass[../main.tex]{subfiles}

\begin{document}

\section{Task 4: Optimal Staging}\label{sec:task4}

\subsection{Staging Model}

Task~4 removes the vertical-climb phase and introduces a single \emph{staging event} at time $t = t_s$. At staging, the first-stage structural mass
\begin{equation}
    m_{s1} = \eta\, Q\, t_s
    \label{eq:ms1}
\end{equation}
is jettisoned instantaneously. Here $Q\,t_s$ is the propellant the first stage burns up to $t_s$ and $\eta = m_s/m_p$ is the \emph{structural coefficient} (inert mass per unit propellant, $\eta = 0.1$): to hold that propellant the stage must carry a tankage/engine mass $\eta\,Q\,t_s$. The propellant is fully consumed---the mass falls as $m(t) = 1 - Q t$---and only the inert structure is dropped, so the vehicle mass undergoes a step:
\begin{equation}
    m^+(t_s) = m^-(t_s) - m_{s1} = 1 - Q\,t_s - \eta\, Q\, t_s = 1 - Q\,t_s(1 + \eta).
\end{equation}
For $m^+(t_s) > 0$ one requires $t_s < [Q(1+\eta)]^{-1}$, which provides an upper bound on the admissible staging time. The costate vector is continuous across the staging event.

The net payload delivered to orbit is
\begin{equation}
    m_\mathrm{pay} = m_f - m_{s2}, \qquad m_{s2} = \eta\, Q\,(t_f - t_s),
    \label{eq:payload}
\end{equation}
where $m_{s2}$ is the second-stage structural mass (again $\eta$ times the propellant that stage burns). Summing the budget over both stages---a total propellant $Q\,t_f$, its structure $\eta\,Q\,t_f$, and the payload, all drawn from $m_0 = 1$---collapses to the compact relation
\begin{equation}
    m_\mathrm{pay} = 1 - (1+\eta)\,Q\,t_f,
    \label{eq:pay_budget}
\end{equation}
which depends only on the \emph{total} burn time $t_f$ (equivalently, the total propellant). Staging raises the payload precisely because shedding the spent first-stage structure lets the vehicle reach orbit on a smaller total propellant load, i.e.\ a smaller $t_f$: here $t_f$ drops from $0.444$ (single stage) to $0.424$ (two-stage), lifting $m_\mathrm{pay}$ from $0.024$ to $0.068$.

\subsection{Optimization Procedure}

For a fixed staging time $t_s$, the two burn arcs share the same costate evolution (the bilinear-tangent law is unaffected by staging), and the BVP reduces to the same structure as Task~1 with updated initial mass after staging:
\begin{enumerate}
    \item Integrate the augmented state from $t = 0$ to $t = t_s$ (Stage~1 burn),
    \item Apply the staging mass jump \eqref{eq:ms1},
    \item Integrate from $t_s$ to $t_f$ (Stage~2 burn),
    \item Evaluate the residual \eqref{eq:residual1} at $t_f$.
\end{enumerate}
The shooting unknowns $\boldsymbol{\zeta}$ (Eq.~\ref{eq:unknowns}) are identical to Task~1.

The optimal staging time $t_s^*$ is found by an \textbf{outer parametric sweep}: $t_s$ is varied over a dense grid of 50 points, the inner BVP is solved at each point (warm-started from the previous converged solution via continuation), and the payload \eqref{eq:payload} is maximized over the converged solutions. The starting guess for the inner BVP at the first grid point uses the single-stage solution from Task~1.

\subsection{Results}

Figure~\ref{fig:task4_pay} shows the payload as a function of staging time $t_s$ for $Q = 2$, $y_f = 0.04$. The curve has a well-defined maximum; the single-stage baseline (dashed red) lies significantly below the two-stage peak.

\begin{figure}[H]
    \centering
    \includegraphics[width=0.78\textwidth]{task4_task_4_payload_vs_staging_time}
    \caption{Task 4: payload mass vs.\ staging time $t_s$. Red dashed line: single-stage reference. Red star: optimal staging point.}
    \label{fig:task4_pay}
\end{figure}

Table~\ref{tab:task4} summarises the numerical results at $Q = 2$, $y_f = 0.04$.

\begin{table}[H]
\centering
\caption{Task 4 key numerical results ($Q = 2$, $y_f = 0.04$).}
\label{tab:task4}
\begin{tabular}{lll}
\toprule
Quantity                  & Value  & Notes \\
\midrule
Optimal staging time $t_s^*$ & 0.337 & nondim ($\approx 268$~s) \\
Single-stage payload      & 0.0241 & nondim \\
Two-stage payload         & 0.0680 & nondim \\
Payload gain              & $+182\%$ & vs.\ single-stage \\
\bottomrule
\end{tabular}
\end{table}

Figures~\ref{fig:task4_traj} and~\ref{fig:task4_mass} show the optimal two-stage trajectory and mass profile. The discontinuous mass drop at $t_s^*$ is visible in the profile; the trajectory is continuous (position and velocity are unaffected by staging, only mass changes).

\begin{figure}[H]
    \centering
    \begin{subfigure}[b]{0.48\textwidth}
        \centering
        \includegraphics[width=\textwidth]{task4_task_4_optimal_trajectory}
        \caption{Optimal two-stage trajectory: Stage~1 (blue), Stage~2 (red), staging point (circle).}
        \label{fig:task4_traj}
    \end{subfigure}
    \hfill
    \begin{subfigure}[b]{0.48\textwidth}
        \centering
        \includegraphics[width=\textwidth]{task4_task_4_mass_profile}
        \caption{Mass profile with the instantaneous jettison at $t_s^*$.}
        \label{fig:task4_mass}
    \end{subfigure}
    \caption{Task 4: optimal two-stage solution for $Q = 2$, $y_f = 0.04$.}
    \label{fig:task4}
\end{figure}

The large relative gain is a consequence of the model simplifications (constant $Q$, no drag, identical $I_{sp}$ for both stages). Real launchers achieve much smaller but still substantial staging benefits; the key mechanism---jettisoning inert structural mass before it has to be accelerated to orbital velocity---is the same.

\end{document}
